problem:  
Let \( M^5 \) be a closed, simply connected Riemannian manifold with **nonnegative sectional curvature** and an effective, isometric **maximal torus action** \(T^3 \curvearrowright M\).  The orbit space \(M/T^3\) is a compact 2-dimensional Alexandrov space with boundary, and each boundary edge corresponds to an isotropy circle subgroup \(S^1_{u_i} \subset T^3\) determined by a primitive integer vector \(u_i \in \mathbb{Z}^3\).  Consecutive edges meeting at isolated  \(T^2\)-fixed orbits must satisfy a **smoothness (unimodularity) condition**: any pair \((u_i, u_{i+1})\) extends to a \(\mathbb{Z}\)-basis of \(\mathbb{Z}^3\).  

Topologically, such data (together with cyclic order) determine the 5‑manifold \(M\) up to equivariant diffeomorphism (Oh 1982).  However, whether a given topologically valid isotropy pattern can actually occur under a nonnegatively curved metric remains open.

**Open problem:**  
Among all topologically admissible weighted boundary data \((u_1,\dots,u_k)\) satisfying the unimodularity conditions from the classification of smooth \(T^3\)-actions on simply connected 5‑manifolds, determine precisely which configurations can be realized by manifolds \(M^5\) admitting nonnegative curvature and which cannot.  

Equivalently: does nonnegative curvature impose *additional combinatorial or number‑theoretic restrictions* on the isotropy slope data beyond those required for topological smoothness?  

Why is it a "good" problem:  
This formulation cleanly separates two layers of structure—**topological admissibility** versus **geometric realizability**.  The first is essentially solved (the slope data that can topologically occur), whereas the second, involving curvature, remains an open frontier.  

It is “good” because:  
- **Profound insight:** It asks whether curvature rigidity is reflected at the **discrete level of lattice isotropy data**, translating geometric smoothness into arithmetic constraints—a potential bridge between Riemannian geometry and lattice combinatorics.  
- **Bridge between fields:** The problem connects **nonnegative curvature geometry**, **equivariant topology**, and **integer lattice theory**. Any progress requires understanding how Alexandrov convexity and metric splitting phenomena restrict the combinatorial possibilities allowed by smooth topology.  
- **Extreme simplicity:** The entire question can be expressed as “Which integral polygons in \(\mathbb{Z}^3\) with unimodular consecutive edges can arise from nonnegatively curved 5‑manifolds?”—visibly simple yet profoundly tied to open rigidity questions in geometry and classification.